\chapter{Implementation Verification} \label{app:verification}

This appendix consolidates the three verification experiments that confirm
BIP-specified behavior holds in the reference implementations. These are
implementation correctness checks, not security findings---the outcomes are
predetermined by the opcode semantics and cryptographic constructions.

% ======================================================================
\section{CAT+CSFS Hot Key Theft Resistance} \label{ase:hot_key}

Experiment~M verifies that the dual CSFS+CHECKSIG verification in the CAT+CSFS vault
prevents hot key holders from redirecting withdrawal outputs. The test modifies the
withdrawal destination in the trigger transaction and attempts to broadcast it. The
dual verification rejects the modified transaction because changing the output changes
the sighash, and the same Schnorr signature cannot satisfy both
\texttt{OP\_CHECKSIGFROMSTACK} (against the stack-assembled preimage containing the
original \texttt{sha\_single\_output}) and \texttt{OP\_CHECKSIG} (against the real
transaction sighash with the modified output).

\textbf{Result:} All output redirection attempts are rejected at the consensus level.
The hot key holder can only trigger to the pre-committed vault-loop output (griefing),
not redirect funds. This confirms the dual-verification binding described in
BIP-347~\cite{BIP347} and BIP-348~\cite{BIP348}.

% ======================================================================
\section{CAT+CSFS Witness Manipulation Resistance} \label{ase:witness}

Experiment~N tests three independent tampering vectors against the CAT+CSFS
sighash preimage:

\begin{enumerate}
\item \textbf{nVersion tampering.} The witness-provided transaction version field
  is modified. The assembled preimage no longer matches the real transaction sighash;
  \texttt{OP\_CHECKSIG} rejects the signature.

\item \textbf{Codeseparator position tampering.} The \texttt{codesep\_pos} field in
  the witness is altered. The CSFS check against the assembled preimage passes (the
  signature was computed for the original preimage), but the CHECKSIG check fails
  because the real sighash uses the actual \texttt{codesep\_pos}.

\item \textbf{Hash type tampering.} The sighash type byte is modified from
  \texttt{SIGHASH\_SINGLE\allowbreak|ANYONECANPAY} (0x83) to \texttt{SIGHASH\_ALL} (0x01).
  Both checks reject because the signature was computed for the original hash type.
\end{enumerate}

\textbf{Result:} All three tampering vectors are rejected. The dual-verification
binding ensures that any byte-level modification to the witness-provided preimage
fields is detected by at least one of the two signature checks. This follows from
SHA-256 collision resistance: the attacker would need to find a preimage that hashes
to the same sighash under a different transaction structure.

The 520-byte \texttt{OP\_CAT} consensus limit constrains the assembled preimage to
520~bytes. The BIP-342 sighash preimage (including the TapSighash tag) is 227~bytes,
leaving 293~bytes of headroom.

% ======================================================================
\section{CCV Mode Bypass} \label{ase:mode_bypass}

Experiment~I verifies that undefined CCV mode values trigger \texttt{OP\_SUCCESS},
bypassing all covenant enforcement. Five undefined mode values are tested:
$\{3, 4, 7, 128, 255\}$. For each, a production-shaped vault taptree is constructed
with the test mode value on the recover leaf. The attacker constructs a recovery
transaction redirecting funds to an arbitrary address with no signature.

\textbf{Result:} All five undefined modes produce complete covenant bypass. The
script evaluates to true unconditionally (\texttt{OP\_SUCCESS} semantics), and funds
are transferred to the attacker-controlled address at 110\vB{} per spend with zero
authentication.

This behavior is specified in BIP-443~\cite{BIP443}: ``Any other value of the mode
makes the opcode succeed validation immediately.'' The design intent is
forward-compatible soft-fork extensibility---undefined modes are reserved for future
opcode semantics. The risk is to developers writing raw CCV scripts who use a mode
value outside the defined set $\{-1, 0, 1, 2\}$. The pymatt library mitigates this
by exposing named constants (\texttt{CCV\_FLAG\_CHECK\_INPUT}, etc.) that prevent
accidental use of undefined values.
